\justifying
\NSFCBodyText

\subsubsection{已具备的实验条件}

\subsubsubsection{计算条件}

申请人已确认可稳定使用的计算资源包括 NVIDIA A40 图形处理器 8 块与 NVIDIA GeForce RTX 3090 Ti 图形处理器 6 块。本项目采用参数高效适配策略控制可训练参数规模，不依赖超大规模集群；上述资源可支撑中等规模视觉语言模型的领域适配、三臂受控对照，以及每组实验不少于 3 个随机种子的重复运行与多组消融的排队执行。

\subsubsubsection{软件与工程条件}

申请人具备 Linux、Python、主流深度学习框架、Git 与 \LaTeX{} 工具链的完整使用经验；在软件系统研发、数据治理与科研工程化方面的积累，可直接用于本项目的数据审计流程、标注工具搭建、实验日志管理与轻量检索演示开发。

\subsubsubsection{数据条件}

本项目使用公开遥感图文数据。RSITMD\cite{yuan2022amfmn} 与 RSICD\cite{lu2018rsicd} 的作者仓库均提供影像、标注文件与解析程序的公开获取入口，相关发布页面已完成核验；必要且许可条件满足时，可使用 SkyScript\cite{wang2024skyscript} 作小规模补充验证。本项目不依赖需要付费采购的商业数据，也不依赖大规模实地采集。

\subsubsubsection{团队与依托单位条件}

申请人为新疆大学副教授、硕士研究生导师，创建并负责 ALPHA Group 研究团队，具备指导研究生开展系统性研究的基础。依托单位新疆大学在计算机与人工智能方向具备相应的学科平台与研究生培养条件，申请人所在二级单位为\textbf{【待补 F-MAP-04：申请人所在二级单位的正式名称与人事口径】}。

\subsubsection{尚缺少的实验条件及拟解决的途径}

\subsubsubsection{数据的实际可用性尚未完成验证}

如前所述，公开数据集目前仅完成发布页面与获取入口层面的核验，其空间关系覆盖率、原始素材授权条款与可再分发范围均属\textbf{【待补 F-MAP-02：公开数据的实际获取结果、许可条款与关系证据可用性结论】}。

\textbf{拟解决途径}：将数据获取、完整性核对与许可确认作为第一年的首批任务，与数据审计同步完成。若某一数据集的许可不支持派生标注公开，则相应成果改以数据索引、处理脚本与版本说明的形式交付；若关系证据覆盖程度不足以支撑三态判定，则按既定退出路径调整研究内容一的定位。该风险已在研究方案中预置了应对安排。

\subsubsubsection{关系证据的人工标注力量需要组织}

内容一需要建立经人工核验的关系挑战子集，标注量较大且要求双人独立标注与一致性检验，目前尚未形成固定的标注队伍与平台。

\textbf{拟解决途径}：依托 ALPHA Group 的研究生力量组织标注小组，先在小规模样本上完成规范试点与一致性评估再扩大规模；控制关系词表规模以降低标注难度；自行搭建轻量标注与仲裁工具，不采购商业标注服务。

\subsubsubsection{补充算力的落实情况}

除申请人已确认可稳定使用的上述设备外，学院层面存在可申请的公共算力资源，但\textbf{固定账号、设备型号与使用额度尚未落实}，因此本申请仅将其作为补充条件列出，不作为稳定可用的实验条件计入。

\textbf{拟解决途径}：本项目的实验设计以已确认可用的设备为基准进行规划，不依赖补充算力即可完成既定研究任务；若后续获得学院公共算力的正式配额，将用于扩大消融实验的覆盖范围与随机种子数量，属于增量而非必要条件。

\subsubsection{利用研究基地的计划与落实情况}

关于利用国家实验室、全国重点实验室和部门重点实验室等研究基地的计划与落实情况，当前状态为\textbf{【待补 F-MAP-05：拟利用的研究基地名称、使用计划与落实情况；如确无相关安排，则据实说明无】}。

这里需要交代清楚：本项目的研究内容基于公开遥感图文数据与申请人已确认可用的计算条件开展，\textbf{不以取得研究基地支持为前提}。若后续落实相关研究基地的使用安排，将主要用于扩大验证规模与加强成果的领域对接，可以做得更充分，但不影响既定任务的完成。
